% ============================================================
% II. RELATED WORK
% ============================================================
\section{Related Work}
\label{sec:related}

% ---- 1. URDF and Robot Description Tooling ----
The Unified Robot Description Format (URDF)~\cite{urdf_doc} is the dominant XML-based representation for kinematic trees in the ROS ecosystem.
Standard URDF (as implemented by urdfdom~\cite{urdfdom}) supports spheres, cylinders, boxes, and meshes but lacks a native capsule element, requiring capsule-based approximations to be decomposed into cylinder-plus-sphere primitives. The Tesseract robotics framework extends URDF parsing with custom \texttt{<capsule>} and \texttt{<convex\_mesh>} tags, representing an alternative parser-extension approach to our standard-primitive decomposition strategy. Xacro~\cite{xacro} provides macro preprocessing but the output remains standard URDF. Alternative formats such as MJCF~\cite{todorov2012mujoco} and SDFormat~\cite{sdformat} offer richer geometry primitives, yet URDF remains the de facto standard for ROS-compatible manipulators.

% ---- 2. Convex Decomposition and Hull Approximation -%
Convex decomposition for collision geometry is a well-studied problem.
Approximate convex decomposition (ACD)~\cite{lien2008acd} partitions non-convex meshes into near-convex components.
The V-HACD algorithm~\cite{mamou2016vhacd} introduced hierarchical decomposition with concavity thresholds and is widely used in robotics and game engines.
CGAL~\cite{cgal} provides robust convex hull computation via the Quickhull algorithm~\cite{barber1996quickhull}, which forms the basis of the convex mode in our toolchain.
Libigl~\cite{jacobson2018libigl} offers a unified geometry-processing interface to these algorithms.
While convex hulls provide tight mesh-based approximations suitable for most simulators, they are still meshes and do not admit the closed-form distance evaluation of analytic primitives.

% ---- 3. Sphere-Tree Approximation -%
Sphere trees are a classic collision geometry representation that approximates shapes with hierarchies of bounding spheres.
The medial-axis sphere tree~\cite{bradshaw2002sphere, mlund_spheretree} constructs an inner approximation whose spheres lie within the original mesh, providing conservative collision estimates.
Spherical approximations are particularly attractive for GPU-based collision pipelines and for signed-distance computations via sphere-sums~\cite{koptev2023neural}.
Our sphere-tree mode uses the vendored \texttt{sphere\_tree} library~\cite{mlund_spheretree} to construct medial-axis trees with configurable density, outputting both a URDF with sphere primitives and a JSON sidecar.
More recently, Foam~\cite{coumar2025foam} introduced an adaptive medial-axis approximation method specialized for spherical approximation of robot URDF geometries, sharing our goal of automating primitive generation for robot collision models. Unlike Foam's spherical-only output, our toolchain additionally produces capsule and convex approximations with analytic parameter export through a JSON sidecar.

% ---- 4. Capsule and Swept-Sphere Primitives -%
Capsule (swept-sphere) primitives have received less attention in automatic fitting pipelines despite their utility.
Wu et al.~\cite{wu2018capsule} introduced a cross-section-based method for fitting capsules to triangular meshes by slicing along a principal axis, fitting circles to each cross-section contour, and chaining circles into capsules---the approach that forms the conceptual foundation of our method.
The capsule distance function (point-to-line-segment minus radius) is a special case of the more general swept-sphere family~\cite{larsen2000ssv}, which includes rectangle-swept spheres (RSS) and line-swept spheres (LSS).
The gProximity library~\cite{lauterbach2010gproximity} exploits the swept-sphere representation for GPU-accelerated hierarchical collision and distance queries, demonstrating the computational advantages of analytic primitives over mesh-based collision. In the robotics literature, capsule collision primitives have been used in several simulation and planning frameworks~\cite{drake}, but the generation of these primitives from existing robot URDF meshes has remained a manual process.
Our work addresses this gap by providing an end-to-end pipeline that integrates PCA sectioning, adaptive circle fitting with circle-outside-area (COA) thresholding, and tightness-driven refinement into a URDF-compatible tool.

% ---- 5. Tooling Gap -%
While each component---convex decomposition, sphere-tree generation, and capsule fitting---has been addressed individually, prior work does not combine them into a single URDF-to-URDF pipeline with sidecar analytic parameter export, built-in validation metrics, and preset-based reproducible tuning. The MathWorks capsuleApproximation function (R2022b) provides capsule fitting for rigid body trees but is closed-source and not URDF-native. Our work fills this specific integration gap for the open-source ROS URDF ecosystem.
